\documentclass[12pt]{article}

\usepackage[margin=1.12in]{geometry}
\usepackage{amsmath,amssymb,amsthm,mathtools}
\usepackage{array}
\usepackage{booktabs}
\usepackage{longtable}
\usepackage{hyperref}

\newtheorem{theorem}{Theorem}
\newtheorem{corollary}{Corollary}
\theoremstyle{definition}
\newtheorem{definition}{Definition}
\theoremstyle{remark}
\newtheorem{remark}{Remark}

\newcommand{\A}{\mathcal A}
\newcommand{\C}{\mathbb C}
\newcommand{\D}{\mathcal D}
\newcommand{\F}{\mathcal F}
\newcommand{\M}{\mathcal M}
\newcommand{\N}{\mathcal N}
\newcommand{\W}{\mathcal W}
\newcommand{\Z}{\mathcal Z}

\setlength{\emergencystretch}{3em}
\hypersetup{
  pdftitle={A Finite-Center No-Go Boundary for Quotient-Visible Reconstruction},
  pdfsubject={Operator algebras, finite centers, sufficiency, and AQFT reconstruction},
  pdfkeywords={von Neumann algebras, finite centers, Blackwell sufficiency, Petz sufficiency, AQFT, binary rank}
}

\title{A Finite-Center No-Go Boundary for Quotient-Visible Reconstruction}
\author{}
\date{Live review note of June 20, 2026}

\begin{document}
\maketitle

\begin{abstract}
This note states the smallest Team 1 residue in standard language.  Let \(A\)
be the represented von Neumann algebra of supplied observable data, and let
\(\widehat A_q=A\oplus (N\overline\otimes \ell^\infty(F_q))\) be a finite
central extension with quotient map \(\pi_q:\widehat A_q\to A\).  Any record
formed only after applying \(\pi_q\) is constant on the hidden finite central
fiber.  Therefore no rule whose input is only that quotient-visible record can
determine targets that vary on the fiber: hidden central support, ambient
faithfulness, finite central cardinality, the predicate \(q=2\), unquotiented
tensor coordinates, omitted sector-character fibers, or absolute finite-center
entropy calibration.

The proof is elementary: normal quotients of von Neumann algebras are
central-projection quotients, and reconstruction from a record is possible only
for functions constant on its fibers.  The result should therefore be treated
as demotion-ready.  In statistical language it is a Blackwell/Le Cam/Petz/
Buscemi sufficiency boundary for a declared experiment, channel, subalgebra, or
model.  The possible contribution is not new operator algebra; it is the
black-hole/AQFT warning that recovery, sector, modular, locally covariant, or
entropy records reconstruct only what their hypotheses supply.  Binary finite
rank follows only after a separator, physical-exclusion/maximality premise,
finite-alphabet law, or direct measurement channel is added.
\end{abstract}

\paragraph{Status.}
This is live non-frozen review metadata.  It does not modify either frozen Team
1 paper, Lean archive, public mirror manifest, or SHA-256 package.  Its purpose
is external classification: known theorem, false statement, missing hypothesis,
too broad physically, or useful expository corollary.

\section{Question}

Standard reconstruction theorems reconstruct supplied structure: a represented
algebra, a functor, a sector category, a code algebra, a channel, a split
inclusion, a trace-normalized Type II algebra, or a measured spectrum.  Team 1's
remaining question is the converse:

\begin{quote}
If the displayed record is already quotient-visible with respect to a hidden
finite central completion, what further premise is required before one may
reconstruct hidden finite central rank, ambient faithfulness, unquotiented
coordinates, or the binary predicate \(q=2\)?
\end{quote}

The word ``reconstruct'' means ``is a function of the displayed data on the
declared comparison class.''  There is no obstruction if the model supplies the
full finite center, the atomic measurement, or \(q\) itself as input.  The
obstruction is to inferring those targets from records that have already
forgotten finite central directions.

\section{Same-Data Finite-Center Extensions}

\begin{definition}[quotient-visible record]
Let \(A\) and \(N\) be von Neumann algebras, let \(F_q\) be a finite set with
\(|F_q|=q\), and set
\[
  \widehat A_q=A\oplus \left(N\overline\otimes \ell^\infty(F_q)\right).
\]
Let
\[
  \pi_q:\widehat A_q\to A,\qquad
  \pi_q(a,x)=a
\]
be the normal quotient map.  A record \(D_q\) is quotient-visible if every
state, effect, channel, sector label, recovery map, modular datum, tensor
coordinate, spectroscopy datum, or entropy quantity used to form \(D_q\)
factors through \(\pi_q\).  Equivalently, \(D_q\) is constant on the kernel
fiber \(0\oplus (N\overline\otimes \ell^\infty(F_q))\).
\end{definition}

\begin{theorem}[same-data finite-center obstruction]
\label{thm:same-data}
Let \(\mathcal C\) be a comparison class containing at least two quotient maps
\(\pi_q:\widehat A_q\to A\) as above, or hidden refinements of one such finite
central summand, with quotient-visible records \(D_q\) having the same
displayed value after applying \(\pi_q\).  Then no reconstruction rule whose
input is only that displayed record can determine any target that is not
constant on the comparison fiber.  In particular it cannot determine:
\begin{enumerate}
  \item whether the represented algebra \(A\) was ambient-faithful;
  \item the nonzero finite central support in the hidden summand;
  \item the hidden finite central cardinality \(q\);
  \item the predicate \(q=2\);
  \item unquotiented tensor coordinates in the hidden finite central summand;
  \item omitted central-character sector fibers; or
  \item an absolute finite-center entropy calibration attached to that summand.
\end{enumerate}
Those targets become determinate only after adding a datum that separates,
excludes, identifies as null or gauge, calibrates, or directly measures the
hidden finite central directions.
\end{theorem}

\begin{proof}
The kernel of \(\pi_q\) is the normal central summand
\[
  0\oplus \left(N\overline\otimes \ell^\infty(F_q)\right).
\]
More generally, ultraweakly closed two-sided ideals in von Neumann algebras are
central-projection summands, so ordinary normal quotients forget precisely a
central summand.  Since \(D_q\) factors through \(\pi_q\), all hidden choices in
that summand give the same displayed record.  But the targets listed in the
theorem vary over same-record completions: take \(q=2\) and \(q=3\), vary the
hidden central support, refine an omitted central-character fiber, change an
unrepresented tensor coordinate, or change an additive finite-center entropy
normalization while preserving the quotient record.  A function of the
displayed record must assign the same value to all same-record completions, so
it cannot equal a target that differs between them.  The only exits are
premises that remove the comparison fiber: a separating readout, a maximality
or physical-exclusion theorem, a null/gauge identification, a calibration
rule, a finite-alphabet law, or a direct measurement channel.
\end{proof}

\begin{corollary}[binary rank is not a quotient-visible consequence]
\label{cor:binary}
In the comparison class of Theorem~\ref{thm:same-data}, the predicate \(q=2\)
is reconstructible from quotient-visible records only if the records are
injective on the binary-versus-nonbinary fiber, or if an added physical premise
collapses that fiber.  Entropy, recovery, sector, modular, or locally
covariant data that are identical for \(q=2\) and \(q=3\) do not force binary
rank.
\end{corollary}

\section{A Three-Atom Counterexample to Binary Inference}

The obstruction can be seen without any infinite-dimensional machinery.  Let
\[
  Z_3=\ell^\infty(\{0,1,2\}),\qquad
  V_3=Z_3/\C 1 .
\]
Define a binary pooled readout
\[
  R_3([x_0,x_1,x_2])=x_0-\frac{x_1+x_2}{2}.
\]
Then \(R_3\) has rank one, but \(V_3\) has dimension two and
\[
  [0,1,-1]\in \ker R_3 .
\]
Thus the readout sees exactly one accounted contrast: atom \(0\) versus the
pooled pair \(\{1,2\}\).  It does not see the contrast separating atoms \(1\)
and \(2\).

Equivalently, at the level of central probability vectors,
\[
  (p_0,p_1,p_2)\longmapsto (p_0,p_1+p_2)
\]
has exactly the same displayed binary statistics as the genuine two-atom
center \(\ell^\infty(\{0,1\})\).  Every decision rule, entropy record, recovery
record, sector label, modular datum, or spectroscopy record that depends only
on the pooled binary data gives the same answer for the two-atom model and the
three-atom model with hidden split \(1|2\).  The accounted contrast rank is
one in both cases, but \(q=2\) in one model and \(q=3\) in the other.

This is the finite counterexample behind Corollary~\ref{cor:binary}.  A binary
quotient is not a binary center.  To conclude \(q=2\), one needs not only a
one-dimensional accounted readout but also a theorem or measurement saying the
pooled quotient is the full physical finite center.

\section{Positive Exits}

The theorem is a no-go only relative to quotient-visible data.  It is killed,
demoted, or made inapplicable by any of the following positive hypotheses:

\begin{longtable}{>{\raggedright\arraybackslash}p{0.35\linewidth}>{\raggedright\arraybackslash}p{0.53\linewidth}}
\toprule
Positive input & Effect \\
\midrule
Faithful finite-center readout & Separates the finite-center contrast space;
then the accounted rank is recoverable. \\
\addlinespace
Physical maximality or source faithfulness & Says the represented algebra is
the full physical algebra, not a quotient of a hidden central extension. \\
\addlinespace
Split, type-I, or tensor-independence hypothesis & Supplies spatial tensor
coordinates instead of asking generated-algebra data to reconstruct them. \\
\addlinespace
Full sector category or full code algebra & Makes DHR/DR or OAQEC/JLMS
reconstruction apply to the intended object, rather than to an omitted quotient. \\
\addlinespace
Trace or area calibration & Fixes the finite-center entropy normalization. \\
\addlinespace
Finite-alphabet law or direct spectroscopy & Measures the relevant finite
fiber and can make \(q=2\) a genuine output rather than an inference from
quotient-visible entropy. \\
\bottomrule
\end{longtable}

\section{Demotion Boundary}

The mathematical core of Theorem~\ref{thm:same-data} is not a Clay- or
Nobel-level new theorem.  It is the central-projection quotient theorem plus
the elementary fact that functions of a record are constant on record fibers.
Blackwell comparison of experiments and Le Cam deficiency theory say the same
thing in decision-theoretic language \cite{Blackwell1951,Blackwell1953,
LeCam1964,LeCam1986}.  Petz sufficiency and later quantum comparison theorems
say the noncommutative version relative to a declared subalgebra, channel, or
statistical model \cite{Petz1986,Petz1988,Shmaya2005,Buscemi2012}.

The possible Team 1 value is therefore a precise black-hole/AQFT claim
boundary:

\begin{quote}
Recovery, sector, modular, locally covariant, split, spectroscopy, or entropy
records reconstruct the object supplied by their hypotheses.  They do not, by
themselves, certify maximality against hidden finite central extensions or
derive binary rank.  To force \(q\) or \(q=2\), the finite-center separator,
exclusion, calibration, or direct measurement premise must be named.
\end{quote}

\section{Reviewer Question}

Is this note merely a known sufficiency/central-quotient corollary under
another name, or is the black-hole/AQFT translation a useful claim boundary?
The desired external response is one of:

\begin{itemize}
  \item \emph{known}: cite the exact theorem and demote the note;
  \item \emph{false}: give the counterexample or hidden assumption;
  \item \emph{missing hypothesis}: name the positive premise that makes the
    target reconstructible;
  \item \emph{too broad}: name the physical axiom always assumed in the
    intended setting;
  \item \emph{useful}: say the finite-center no-go is worth writing as an
    expository reconstruction boundary.
\end{itemize}

\begin{thebibliography}{99}

\bibitem{Blackwell1951}
D. Blackwell,
Comparison of experiments,
in \emph{Proceedings of the Second Berkeley Symposium on Mathematical
Statistics and Probability}, University of California Press, 1951, 93--102.

\bibitem{Blackwell1953}
D. Blackwell,
Equivalent comparisons of experiments,
\emph{Annals of Mathematical Statistics} \textbf{24} (1953), 265--272.

\bibitem{LeCam1964}
L. Le Cam,
Sufficiency and approximate sufficiency,
\emph{Annals of Mathematical Statistics} \textbf{35} (1964), 1419--1455.

\bibitem{LeCam1986}
L. Le Cam,
\emph{Asymptotic Methods in Statistical Decision Theory},
Springer, 1986.

\bibitem{Petz1986}
D. Petz,
Sufficient subalgebras and the relative entropy of states of a von Neumann
algebra,
\emph{Communications in Mathematical Physics} \textbf{105} (1986), 123--131.

\bibitem{Petz1988}
D. Petz,
Sufficiency of channels over von Neumann algebras,
\emph{Quarterly Journal of Mathematics} \textbf{39} (1988), 97--108.

\bibitem{Shmaya2005}
E. Shmaya,
Comparison of information structures and completely positive maps,
\emph{Journal of Physics A: Mathematical and General} \textbf{38} (2005),
9717--9727.

\bibitem{Buscemi2012}
F. Buscemi,
Comparison of quantum statistical models: equivalent conditions for
sufficiency,
\emph{Communications in Mathematical Physics} \textbf{310} (2012), 625--647.

\end{thebibliography}

\end{document}
